\section{Frequently Asked Questions}
\label{sec:freq-asked-quest}

\subsection{Session types vs.\ protocol types}
\label{sec:what-diff-betw-1}
 A session type can be assigned to a channel end to describe
the entire communication behavior of this channel up to its closing. Thus, a
session type classifies a run-time value, namely a channel end.
A protocol type does not correspond to any run-time value.  It
describes composable behavior for bidirectional communication. A protocol
type becomes part of a session type by binding it to a particular
direction with the $!$ and $?$ operators.

In a sense, a protocol type describes a composable part of a
communication protocol or the difference between two session
types. This view is emphasized by our pervasive use of channel passing
style. 

\subsection{Negated recursion}
\label{sec:what-about-protocol}
It is possible to define an algebraic protocol that negates the polarity of its recursive use.
Such a protocol is fine, though unusual as it changes
direction with every unfolding of the recursion. Here is an example:
\begin{lstlisting}
protocol Flipper = Flipper -Int -Flipper
\end{lstlisting}
A server for \lstinline+Flipper+ can be implemented with a pair of
mutually recursive functions or even with just one function:
\begin{lstlisting}
flipper : !Flipper.EndT -> Unit
flipper c = let c = select Flipper [EndT] c in
            let (x, c) = receive [Int, ?Flipper.EndT] c in
	    match c with {
	      Flipper c -> send [Int, !Flipper.EndT] x c |> flipper }
\end{lstlisting}

\subsection{Mutual recursion}
\label{sec:do-you-support}
Mutually recursive protocols pose no problem. In general, the server
would be implemented by mutually recursive functions.
% In our recoding of the example in
% \cref{sec:what-about-protocol}, a single recursive function is sufficient:
% \begin{lstlisting}
% protocol Flip = Flip -Int Flop
% protocol Flop = Flop Int Flip

% flipflop : !Flip.End -> End
% flipflop c = let c = select Flip [End] c in
%             let (x, c) = receive [Int, !Flop.End] c in
% 	    select Flop [End] c
% 	    |> send [Int, !Flip.End] x []
%             |> flipflop
% \end{lstlisting}
\begin{lstlisting}
protocol Flip = Flip -Int Flop
protocol Flop = Flop Int Flip

flip : !Flip.EndT -> Unit
flip c = select Flip [EndT] c |> receive [Int, !Flop.EndT] |> flop

flop : (Int, !Flop.EndT) -> Unit
flop p = let (x, c) = p in
         select Flop [EndT] c |> send [Int, !Flip.EndT] x |> flip
\end{lstlisting}
\subsection{Duality vs.\ polarities}
\label{sec:what-diff-betw}
The \lstinline+Dual+ operator works on \emph{session types} and it inverts
the direction of communication from the outside. Intuitively, the
\lstinline+Dual+ operator traverses the spine of a session type and
flips all $!$ to $?$ and vice versa. It does \textbf{not} enter the
type of the message ``payload''. For example (with $\typec T$ the
payload type and $\typec S$ the continuation session):
\begin{align*}
  \TDual{(\TIn TS)} &\equiv \TOut T{\TDual S} &   \TDual{(\TOut TS)} &\equiv \TIn T{\TDual S}
\end{align*}
In contrast, the polarity operator \lstinline+-+ works on
\emph{protocol types} and inverts the direction of communication from
the inside. This inversion happens at the boundary where a protocol
type $\proto$ is turned into a payload type for a session:
\begin{align*}
  \TIn {(\TMinus P)}S &\equiv \TOut P S &\TOut {(\TMinus P)}S &\equiv \TIn P S
\end{align*}
The soundness of the type system guarantees that the direction of communication is
finally settled by the time we get to a primitive sending or receiving operation.

Like \lstinline+Dual+, the \lstinline+-+ operator is involutory:
\begin{align*}
  \TDual{(\TDual{S})} &\equiv \typec{S} &   \TMinus{(\TMinus T)} &\equiv \typec{T}
\end{align*}

\subsection{Recursion and duality}
\label{sec:what-about-inter}
\citet{BernardiH14,BernardiH16} discovered a problem in the interaction
between duality and recursive session types, which is exposed most
concisely by the recursive type $\typec{T =\mu X. !X.X}$.
The dual of this type is $\TDual{(\mu X. !X.X)} = \typec{\mu
  X.\TIn{(\mu X. !X.X)}X}$, which can be seen by dualizing the type's
expansion
\begin{equation*}
  \TDual{(\mu X. !X.X)} = \TDual{(!(\mu X.!X.X). (\mu X.!X.X))} \\
                        = \TIn{(\mu X.!X.X)}{\TDual{(\mu X.!X.X)}}
\end{equation*}
Using equations we could define the type $\typec T$ by $\typec{T = \TOut{T}T}$.
Translated to algebraic session types, the protocol underlying $\typec T$ is as follows:
\begin{lstlisting}
protocol X = Mu T X
type T = !X.EndT

selectMu : T -> !T.T
selectMu c = select Mu [EndT] c

dualT : Dual T -> ?X.EndT
dualT c = c

matchMu : Dual T -> ?T.(Dual T)
matchMu c' = match c' with { Mu c' -> c' }
\end{lstlisting}
This behavior matches exactly the outcome for $\typec T$ according to
\citet{BernardiH16}. The function \lstinline+selectMu+ performs the
unfolding of $\typec T$ whereas \lstinline+matchMu+ performs the
unfolding of $\TDual T$. The type of the identity function
\lstinline+dualT+ shows that \lstinline+Dual T = ?X.EndT+. 

\subsection{Efficiency of generic servers}
\label{sec:are-generic-servers}
In \cref{sec:toolb-gener-serv}, we explore an example with generic servers 
to showcase the expressivity of the proposed parametric protocols.
The example has a significant overhead as there are many tagging
operations compared to a manually crafted version.  

The present work is a first step and does not offer a solution
supported by formal theorems. But, let us call a protocol $\proto$
\emph{tagless} if it is non-recursive and declares a  single
constructor tag $\constr\ T_1\dots T_n$. In analogy to Haskell's
\texttt{newtype}, an implementation can reduce the
run-time cost of a tagless protocol to zero by omitting the select and
match operations. 

\subsection{Equirecursive vs. isorecursive session types}
\label{sec:this-session-types}
Most work on session types is using equirecursive types for two reasons. One is historic:
\citet{DBLP:conf/esop/GayH99,DBLP:journals/acta/GayH05} were the first
to study coinductively defined relations on (tail-) recursive session types. They defined the (by now)
standard notion of subtyping and showed that it is
decidable. Subsequent work has built on top of their results. The
other reason is that ``The equirecursive interpretation of a session type guarantees that no
message is required for the unfolding of a recursive type.'' (\citet{DBLP:journals/pacmpl/BalzerP17})

Algebraic protocols conflate the unfolding of the recursion with
the branching, just like algebraic datatypes do.
For protocols that have finite runs, algebraic protocol do not cause extra messages.
As such protocols use recursion guarded by branching, a
message to decide the branching is needed anyway.

The protocol \lstinline+Stream a+ in \cref{sec:param-prot} does cause
extra messages as it (roughly) corresponds to an isorecursive reading of the type $\mu X. !a.X$.

\subsection{Subtyping, type classes, type inference}
\label{sec:can-you-do}

Subtyping is undecidable for context-free session types
\cite{DBLP:journals/toplas/Padovani19}. However, as \algst's session
typing relies mostly on nominal types, there is scope for decidable
subtyping. We have a concrete proposal, but it is not implemented
and its metatheory has not been checked, yet.

The types \lstinline+Send a+ and \lstinline+Recv a+ for sending and
receiving data of type \lstinline+a+ in \cref{sec:transm-param-data}
look like types for a dictionary in a language with type classes or
traits
\cite{DBLP:journals/toplas/HallHJW96,steve18:_rust_progr_languag}. 
%
In future work we plan to explore the addition of implicit arguments as
in Scala \cite{DBLP:journals/pacmpl/OderskyBLBMS18}. 

Type inference is future work.

\subsection{Settling of direction}
\label{sec:settling-direction}
It might be confusing to see a type like $\TIn TS$ and $\TOut TS$, but
with the possibility that the direction changes later on. However,
this behavior is completely predictable. As long as $\typec T$ is a
protocol type variable  $\alpha : \kindp$, the direction can change
because $\alpha$ could be substituted by some protocol $\TMinus {T'}$. 
As soon as $\typec T$ is an ordinary type of kind $\kindt$ (even if
$\typec T$ is a variable), the direction is fixed. For example, in the
type of primitive send
\lstinline+forall(a:M). forall(s:S).  a -> !a.s -> s+, the payload type
\lstinline+a+ has kind $\kindm$, so the direction is fixed. \\

%%% Local Variables:
%%% mode: latex
%%% TeX-master: "main"
%%% End:
